\scanpage{84}
\mspage{80}
% archive sheet number: 83

Il reste à voir l'unicité de $\unsure{\alpha|\varphi}$, qui est immédiate, vu que $\deleted{\unsure{\forall}}\unsure{A|\mathbb{A}|\mathscr{A}}\in\Sigma_E$ est le $\sup$ des $\sigma_E(x)$ qu'il majore (i.e. tels que $x\in A$ \ldots), \deleted{et $\alpha$ doit \unsure{ce}} d'où forcément
\[
\alpha(A)=\operatorname{Sup}_{x\in A}\underbrace{\alpha(\sigma_E(x))}_{\sigma(\unsure{x|*})}\unsure{=|\equiv}\operatorname{Sup}_{x\in A}\sigma(x)
\qquad\text{\unsure{qed}}
\]

\deleted{\uline{Corollaire} Soit $(E,\inserted{\unsure{\leq^{\circ}|\leq},}\mathrel{|{\circ}|})$ un \unsure{disjoncteur} ordonné, $E_0$ une partie très dense de $E$}

\uline{Corollaire\unsure{!|1}} Soit $\deleted{E}(E,\mathrel{|{\circ}|}_E)$ un \unsure{disjoncteur}, $E'$ une partie de $E$, $\mathrel{|{\circ}|}_{E'}$ la relation induite. Conditions équivalentes:

(i) L'image $\deleted{\unsure{E'}}\sigma_E(E')$ de $E'$ dans $\Sigma_E$ est dense, i.e. tout $S\in\Sigma_E$ est le $\sup$ de l'ens. des éléments de $\sigma_E(E')$ qu'il majore.

(ii) \deleted{\unsure{L'}} L'application
\[
S\longrightarrow S\cap E',\quad\Sigma_E\unsure{{-}|\longrightarrow|\longmapsto}\unsure{\mathscr{P}|\mathcal{P}|\mathfrak{P}}(E')
\]
\deleted{applique $\Sigma_E$ dans $\Sigma_{E'}$, et induit une}
\inserted{est \uline{\unsure{injection|surjection}}}
(NB elle induit alors \unsure{$=$|$\simeq$} un iso
\[
\varphi:\Sigma_E\longrightarrow\Sigma_{E'}).
\]
\deleted{\unsure{(NB. ce sera} $=$ un iso)}

(iii) L'application
\[
\begin{aligned}
\psi:\Sigma_{E'}&\longrightarrow\Sigma_E\\
S'&\longmapsto\unsure{\operatorname{supp}^{\circ}_E|\operatorname{supp}^{o}_E}(S')
\end{aligned}
\]
est \unsure{surjective|bijective} (NB ce sera \unsure{$=$|$\simeq$} un iso)

% slanted
\margin{$\longleftarrow$ \unsure{permuter} avec (v) \deleted{\unsure{comme (i bis)}}}
(iv) $E'$ est \uline{dense} dans $E$ pour la relation d'\inserted{pré}ordre canonique de $E$ (induite de celle de $\Sigma_E$ via $\sigma_E$).
